\chapter{2005年辽宁卷}
\section{单选题}
\begin{problem}
复数 \(z = \frac{-1 + \i}{1 + \i} - 1\) 在复平面内， \(z\) 所对应的点在（\quad）

\choices
    {第一象限}
    {第二象限}
    {第三象限}
    {第四象限}
\end{problem}

\begin{problem}
极限 \(\lim f(x)\) 存在是函数 \(f(x)\) 在点 \(x = x_0\) 处连续的（\quad）

\choices
    {充分而不必要的条件}
    {必要而不充分的条件}
    {充要条件}
    {既不充分也不必要的条件}
\end{problem}

\begin{problem}
设袋中有 80 个红球，20 个白球，若从袋中任取 10 个球，则其中恰有 6 个红球的概率为（\quad）

\choices
    {\(\frac{C_{80}^{6}}{C_{100}^{10}}\)}
    {\(\frac{C_{80}^{6}}{C_{100}^{10}}\)}
    {\(\frac{C_{80}^{6}}{C_{100}^{10}}\)}
    {\(\frac{C_{80}^{6}}{C_{100}^{10}}\)}
\end{problem}

\begin{problem}
已知 m、n 是两条不重合的直线， \( \alpha \) 、 \( \beta \) 、 \( \gamma \) 是两个两不重合的平面，给出下列四个命题： \circled{1} 若 \( m \perp \alpha, m \perp \beta \) ，则 \( \alpha \notin \beta; \) \circled{2} 若 \(\alpha \perp \gamma, \beta \perp \alpha\)，则 \(\alpha \notin \beta\); \circled{3} 若 \(m \subset \alpha, n \subset \beta, m \nmid n\)，则 \(\alpha \notin \beta\); \circled{4} 若 m、n 是异面直线， \( m \subset \alpha, m \notin \beta, n \subset \beta, n \notin \alpha \) ，则 \( \alpha \notin \beta \) . 其中真命题是（\quad）

\choices
    {\circled{1}和\circled{2}}
    {\circled{1}和\circled{3}}
    {\circled{3}和\circled{4}}
    {\circled{1}和\circled{4}}
\end{problem}

\begin{problem}
函数 \( y = \ln (x + \sqrt{x^2 + 1}) \) 的反函数是（\quad）

\choices
    {\(y = \frac{\e^x + \e^{-x}}{2}\)}
    {\( y = -\frac{\e^x + \e^{-x}}{2} \)}
    {\( y = \frac{\e^x - \e^{-x}}{2} \)}
    {\( y = -\frac{\e^x - \e^{-x}}{2} \)}
\end{problem}

\begin{problem}
若 \(\log_2\frac{1 + a^2}{1 + a} < 0\)，则 \(a\) 的取值范围是（\quad）

\choices
    {\(\left(\frac{1}{2}, +\infty\right)\)}
    {\((1, +\infty)\)}
    {\(\left(\frac{1}{2}, 1\right)\)}
    {\(\left(0, \frac{1}{2}\right)\)}
\end{problem}

\begin{problem}
在 R 上定义运算 \( \otimes: x \otimes y = x(1 - y) \) . 若不等式 \( (x - a) \otimes (x + a) < 1 \) 对任意实数 x 成立，则（\quad）

\choices
    {\(-1 < a < 1\)}
    {\(0 < a < 2\)}
    {\(-\frac{1}{2} < a < \frac{3}{2}\)}
    {\(-\frac{3}{2} < a < \frac{1}{2}\)}
\end{problem}

\begin{problem}
若钝角三角形三内角的度数成等差数列，且最大边长与最小边长的比值为 m，则 m 的范围是（\quad）

\choices
    {(1,2)}
    {\((2, +\infty)\)}
    {\([3, +\infty)\)}
    {\((3, +\infty)\)}
\end{problem}

\begin{problem}
若直线 \(2x - y + c = 0\) 按向量 \(\overrightarrow{a} = (1, -1)\) 平移后与圆 \(x^{2} + y^{2} = 5\) 相切，则 \(c\) 的值为（\quad）

\choices
    {8 或 \(-2\)}
    {6 或 \(-4\)}
    {4 或 \(-6\)}
    {2 或 \(-8\)}
\end{problem}

\begin{problem}
已知 \( y = f(x) \) 是定义在 \( \R \) 上的单调函数，实数 \( x_{1} \neq x_{2}, \lambda \neq -1 \)，
\[
\alpha = \frac{x_{1} + \lambda x_{2}}{1 + \lambda}, \quad \beta = \frac{x_{2} + \lambda x_{1}}{1 + \lambda}.
\]
若 \( |f(x_{1}) - f(x_{2})| < |f(\alpha) - f(\beta)| \)，则（\quad）

\choices
    {\(\lambda < 0\)}
    {\(\lambda = 0\)}
    {\(0 < \lambda < 1\)}
    {\(\lambda \geqslant 1\)}
\end{problem}

\begin{problem}
已知双曲线的中心在原点，离心率为 \(\sqrt{3}\)，若它的一条准线与抛物线 \(y^{2} = 4x\) 的准线重合，则该双曲线与抛物线 \(y^{2} = 4x\) 的交点到原点的距离是（\quad）

\choices
    {\(2\sqrt{3} +\sqrt{6}\)}
    {\(\sqrt{21}\)}
    {\(18 + 12\sqrt{2}\)}
    {21}
\end{problem}

\begin{problem}
一给定函数 \( y = f(x) \) 的图象在下列图中，并且对任意 \( a_1 \in (0,1) \)，由关系式 \( a_{n+1} = f(a_n) \) 得到的数列 \( \{a_n\} \) 满足 \( a_{n+1} > a_n (n \in \mathbf{N}^*) \)，则该函数的图象是（\quad）

\bitmapfigure[max width=0.78\linewidth,max totalheight=9.2cm]{img/2005/liaoning/q12_options.png}
\end{problem}

\section{填空题}
\begin{problem}
\(\left(x^{k} - 2x^{-k}\right)^{10}\) 的展开式中常数项是
\end{problem}

\begin{problem}
如图，正方体的棱长为 \(1, C, D\) 分别是两条棱的中点，\(A, B, M\) 是顶点，那么点 \(M\) 到截面 \(ABCD\) 的距离是 \fillinblank{} \fillinblank{}.
\end{problem}

\begin{problem}
用 1、2、3、4、5、6、7、8 组成没有重复数字的八位数，要求 1 和 2 相邻，3 与 4 相邻，5 与 6 相邻，而 7 与 8 不相邻，这样的八位数共有 个. (用数字作答)
\end{problem}

\begin{problem}
\(\omega\) 是正实数. 设 \(S_{\omega} = \{\theta | f(x) = \cos [\omega (x + \theta)]\) 是奇函数 \(\}\)，若对每个实数 \(a, S_{\omega} \cap (a, a + 1)\) 的元素不超过 2 个，且有 \(a\) 使 \(S_{\omega} \cap (a, a + 1)\) 含 2 个元素，则 \(\omega\) 的取值范围是 \fillinblank{} \fillinblank{}.
\end{problem}

\begin{problem}
已知三棱锥 \(P - ABC\) 中，\(E, F\) 分别是 \(AC, AB\) 的中点，\(\triangle ABC, \triangle PEF\) 都是正三角形，\(PF \perp AB\). (1) 证明: \(PC \perp\) 平面 \(PAB\); (2) 求二面角 \(P - AB - C\) 的平面角的余弦值； (3) 若点 \(P, A, B, C\) 在一个表面积为 \(12\pi\) 的球面上，求 \(\triangle ABC\) 的边长 \fillinblank{}.
\end{problem}

\begin{problem}
如图，在直线为 1 的圆 O 中，作一关于圆心对称、邻边互相垂直的十字形，其中 \(y > x > 0\). (1) 将十字形的面积表示为 \( \theta \) 的函数; (2) \(\theta\) 为何值时，十字形的面积最大？最大面积是多少？
\end{problem}

\begin{problem}
已知函数 \( f(x) = \frac{x + 3}{x + 1} (x \neq -1) \). 设数列 \( \{a_n\} \) 满足 \( a_1 = 1, a_{n+1} = f(a_n) \)，数列 \( \{b_n\} \) 满足 \( b_n = |a_n - \sqrt{3}| \)，\( S_n = b_1 + b_2 + \cdots + b_n (n \in \mathbf{N}^*) \). (1) 用数学归纳法证明： \( b_{n} \leqslant \frac{(\sqrt{3}-1)^{n}}{2^{n-1}} \) ; (2) 证明： \( S_{n}<\frac{2\sqrt{3}}{3}. \)
\end{problem}

\begin{problem}
某工厂生产甲、乙两种产品，每种产品都是经过第一和第二工序加工而成，两道工序的加工结果相互独立，每道工序的加工结果均有 \(A\)、\(B\) 两个等级. 对每种产品，两道工序的加工结果都为 \(A\) 级时，产品为一等品，其余均为二等品. (1) 已知甲、乙两种产品每一道工序的加工结果为 A 级的概率如表一所示，分别求生产出的甲、乙产品为一等品的概率 \( P_{\text{甲}} \) 、 \( P_{\text{乙}} \) ; 概率/产品 第一工序 第二工序 甲 0.8 0.85 乙 0.75 0.8 表一 (2) 已知一件产品的利润如表二所示，用 \(\zeta\)、\(\eta\) 分别表示一件甲、乙产品的利润，在 (1) 的条件下，求 \(\zeta\)、\(\eta\) 的分布列及 \(E\xi\)、\(E_{\eta}\); 利润/产品 等级 一等 二等 甲 5(万元) 2.5(万元) 乙 2.5(万元) 1.5(万元) 表二 (3) 已知生产一件产品需用的工人数和资金额如表三所示. 该工厂有工人40名，可用资. 金60万元. 设 \(x\)、\(y\) 分别表示生产甲、乙产品的数量，在(2)的条件下，\(x\)、\(y\) 为何值时，\(z = xEz + yE_{\eta}\) 最大\(\perp\) 最大值是多少\(\perp\) 产品项目 工人(名) 资金(万元) 甲 8 8 乙 2 10 表三
\end{problem}

\begin{problem}
已知椭圆 \(\frac{x^2}{a^2} + \frac{y^2}{b^2} = 1 (a > b > 0)\) 的左、右焦点分别是 \(F_1(-c, 0)\)、\(F_2(c, 0)\)、\(Q\) 是椭圆外的动点. 满足 \(\left|\overrightarrow{F_1Q}\right| = 2a\). 点 \(P\) 是线段 \(F_1Q\) 与该椭圆的交点，点 \(T\) 在线段 \(F_2Q\) 上，并且满足 \(\overrightarrow{PT} \cdot \overrightarrow{TF_2} = 0\)，\(\left|\overrightarrow{F_2T}\right| \neq 0\). (1) 设 \(x\) 为点 \(P\) 的横坐标，证明: \(\left|\overrightarrow{F_1P}\right| = a + \frac{y}{x} a\); (2) 求点 T 的轨迹 C 的方程; (3) 试问: 在点 \(T\) 的轨迹 \(C\) 上，是否存在点 \(M\)，使 \(\triangle F_1MP_2\) 的面积 \(S = b^2\). 若存在，求 \(\angle F_1MP_2\) 的正切值; 若不存在，请说明理由 \fillinblank{}.
\end{problem}

\begin{problem}
函数 \( y = f(x) \) 在区间 \((0, +\infty)\) 内可导，导函数 \( f'(x) \) 是减函数，且 \( f'(x) > 0 \). 设 \( x_0 \in (0, +\infty), y = kx + m \) 是曲线 \( y = f(x) \) 在点 \((x_0, f(x_0))\) 处的切线方程，并设函数 \( g(x) = kx + m \). (1) 用 \(x_0\)、\(f(x_0)\)、\(f'(x_0)\) 表示 \(m\); (2) 证明: 当 \(x_0 \in (0, +\infty)\)，\(g(x) \geqslant f(x)\); (3) 若关于 \(x\) 的不等式 \(x^{2} + 1 \geqslant ax + b \geqslant \frac{1}{2} x^{2}\) 在 \([0, +\infty)\) 上恒成立. 其中 \(a, b\) 为实数，求 \(b\) 的取值范围及 \(a\) 与 \(b\) 所满足的关系 \fillinblank{}.
\end{problem}

